%%% ============================================================
%%% First-Coprime-Failure and the Discrete Fejer Kernel:
%%% A Coordinate Translation across Squarefree Bases
%%%
%%% B. R. Sanders, M. Gish (2026)
%%% Target venue: Experimental Mathematics
%%% MSC: 11A41, 11N05, 11A51, 42A16, 11Y70
%%% DOI of bundle: 10.5281/zenodo.18852047
%%%
%%% Builds on J04 (First-G Law, submitted to Integers) as the
%%% foundational localization lemma. The present note records the
%%% companion harmonic identity (the discrete Fejer kernel at
%%% rational frequency) and the elementary fact that its first
%%% zero co-locates with the first-coprime-failure index. The
%%% closed form is a textbook calculation; the contribution is the
%%% packaging plus the exhaustive verification harness.
%%% ============================================================

\documentclass[11pt,reqno]{amsart}

\usepackage{amsmath,amssymb,amsthm,mathtools}
\usepackage[margin=1in]{geometry}
\usepackage[colorlinks=true,linkcolor=blue,citecolor=blue,urlcolor=blue]{hyperref}
\usepackage{microtype}
\usepackage{enumitem}
\usepackage{booktabs}

%% -------- theorem environments --------
\theoremstyle{plain}
\newtheorem{theorem}{Theorem}[section]
\newtheorem{lemma}[theorem]{Lemma}
\newtheorem{proposition}[theorem]{Proposition}
\newtheorem{corollary}[theorem]{Corollary}

\theoremstyle{definition}
\newtheorem{definition}[theorem]{Definition}
\newtheorem{example}[theorem]{Example}

\theoremstyle{remark}
\newtheorem{remark}[theorem]{Remark}

%% -------- macros --------
\DeclareMathOperator{\spf}{spf}
\DeclareMathOperator{\sinc}{sinc}
\newcommand{\Z}{\mathbb{Z}}
\newcommand{\N}{\mathbb{N}}
\newcommand{\R}{\mathbb{R}}
\newcommand{\C}{\mathbb{C}}
\newcommand{\Prm}{\mathbb{P}}
\newcommand{\bigO}{\mathcal{O}}

%% -------- title matter --------
\title[First-Coprime-Failure and the Discrete Fej\'er Kernel]%
{First-Coprime-Failure and the Discrete Fej\'er Kernel:\\
 A Coordinate Translation across Squarefree Bases}

\author{Brayden R.\ Sanders}
\address{7Site LLC, Hot Springs, Arkansas, USA}
\email{brayden@7site.co}

\author{M.\ Gish}
\address{Independent Researcher, Hot Springs, Arkansas, USA}
\email{monica.gish1992@gmail.com}

\date{2026}

\keywords{coprimality partition, smallest prime factor, Fej\'er
kernel, discrete sinc identity, sieve of Eratosthenes,
exponential sums on rational frequencies, experimental number
theory}

\subjclass[2020]{11A41, 11N05, 11A51, 42A16, 11Y70}

\begin{document}

\begin{abstract}
For an integer $b > 1$ with smallest prime factor $p_1 = \spf(b)$,
the First-G Localization Lemma~\cite{J04Sanders} identifies
$k = p_1$ as the smallest alphabet size at which the sieve of
Eratosthenes marks an element of $\{1, \ldots, k\}$ as
non-coprime to $b$. The present note records a companion
harmonic identity, the discrete Fej\'er kernel at rational
frequency, and the elementary fact that its first integer zero
co-locates with the first-coprime-failure index --- a coordinate
translation between the arithmetic gate at $k = p_1$ and the
harmonic zero of $R(k, p_1) = \sin^2(\pi k / p_1) /
(k^2 \sin^2(\pi / p_1))$ at $k = p_1$, both of which solve
``smallest positive $k$ such that $p_1 \mid k$.'' The closed form
is a standard Fej\'er-kernel calculation
\cite{Fejer1900,Apostol1976,IwaniecKowalski2004}; we record it
together with the Eratosthenes synchronization, an
$\omega$-blindness corollary, and a discrete-to-continuum
sinc-squared identity. We verify all four claims across $8$
small primes and $187$ semiprimes: $712$ distinct algebraic
checks, max deviation $1.11 \times 10^{-16}$ (machine
$\epsilon$), zero exceptions; runtime under three minutes on a
2024 consumer laptop. A brief remark relates the universal
mid-period constant $\sinc^2(1/2) = 4/\pi^2$ to the same
constant in Montgomery's pair correlation
$R_2(u) = 1 - \sinc^2(u)$~\cite{Montgomery1973}, a coincidence of
the rectangular-pulse spectral window common to both contexts.
\end{abstract}

\maketitle

%%% ==================================================================
\section*{0. Lens, substrate, and tier discipline}
%%% ==================================================================

\noindent
\textsc{Lens and substrate.}
The arithmetic side of the present paper works on $\Z$ with the
smallest-prime-factor function $\spf : \Z_{> 1} \to \Prm$, which
is structurally orthogonal to the algebraic substrate
$\Zten = \Z / 10\Z$ studied in companion
work~\cite{J04Sanders,FourCore}. The two sides share only the
broader research program's interest in prime-indexed phase
transitions; the present paper's claims are number-theoretic,
not magma-theoretic, and stand independently of the algebraic
substrate.

\medskip\noindent
\textsc{Tier discipline.}
\begin{itemize}\setlength\itemsep{2pt}
\item \emph{KNOWN (Lemma).} The closed form
$R(k, f) = \sin^2(\pi k / f) / (k^2 \sin^2(\pi / f))$ is the
discrete Fej\'er kernel at rational frequency $1/f$, normalized
by alphabet size. It is in Apostol~\cite{Apostol1976} \S11.5,
Iwaniec--Kowalski~\cite{IwaniecKowalski2004} \S1.7, and
Oppenheim--Schafer~\cite{OppenheimSchafer2010} \S3.8 as the
frequency response of an FIR filter with rectangular window.
We attribute the lemma directly to Fej\'er~\cite{Fejer1900}.
\item \emph{PROVEN.} The synchronization (Theorem~\ref{thm:zerowidth})
recording that the First-G event of~\cite{J04Sanders} at
$k = p_1 = \spf(b)$ and the first integer zero of
$R(\cdot, p_1)$ at $k = p_1$ are the same event in two
coordinate systems, both solving ``smallest positive $k$ such
that $p_1 \mid k$.'' The continuum identification
$R(k, f) \to \sinc^2(k/f)$
(Theorem~\ref{thm:continuum}). The $\omega$-blindness statement
that $R(k, 1/p)$ depends only on $k$ and $p$, not on the
modulus $b$ (Corollary~\ref{cor:omegablind}) --- a one-line
consequence of the closed form.
\item \emph{COMPUTED.} The verification harness exhausts
$712$ distinct algebraic checks: $106$ closed-form $(f, k)$
pairs across $8$ primes, $561$ pre-collapse and gate evaluations
across $187$ semiprimes, $42$ ring-structure cells for the
$\omega$-blindness statement, and $3$ continuum-limit witnesses
for Theorem~\ref{thm:continuum}. Runtime under three minutes;
max deviation $1.11 \times 10^{-16}$ (machine $\epsilon$); zero
exceptions. The script
\texttt{verify\_prime\_phase\_transition.py} runs from the
manuscript folder with no external data dependencies beyond
\texttt{numpy} and the \texttt{cmath} module.
\item \emph{STRUCTURAL RHYME.} The constant $\sinc^2(1/2) =
4/\pi^2 = (2/3)/\zeta(2)$ is a textbook consequence of
$\zeta(2) = \pi^2/6$. We record its appearance at the
half-period of the rectangular spectral window, and we note
without claim that the same constant appears in Montgomery's
pair correlation $R_2(u) = 1 - \sinc^2(u)$ at $u = 1/2$, a
coincidence reflecting the universality of the rectangular
spectral window common to both contexts~\cite{Montgomery1973,
OppenheimSchafer2010}.
\item \emph{OPEN.} What forces the rectangular-pulse alphabet
$\{1, \ldots, k\}$ as the natural object of study (rather than
some smoothed-window analogue) is not addressed here. Whether
arithmetic counterparts of the GUE pair-correlation density
admit closed-form approximations in some scaling regime is
left as an open question (Remark~\ref{rem:open-bridge}).
\end{itemize}

%%% ==================================================================
\section{Introduction}
\label{sec:intro}
%%% ==================================================================

\subsection*{Setup.}

The First-G Localization Lemma of~\cite{J04Sanders} is a clean
foundational lemma about the sieve of Eratosthenes: for an
integer $b > 1$ with smallest prime factor $p_1$, the coprimality
partition
\[
  C_k(b) = \{x \in \{1, \ldots, k\} : \gcd(x, b) = 1\},
  \qquad
  G_k(b) = \{1, \ldots, k\} \setminus C_k(b)
\]
satisfies $|G_k(b)| = 0$ for every $k < p_1$ and
$G_{p_1}(b) = \{p_1\}$, so the alphabet $\{1, \ldots, k\}$ first
contains an element non-coprime to $b$ at exactly $k = p_1$. The
\emph{stability window} of $b$ is the interval
$\{1, \ldots, p_1 - 1\}$ on which $|G_k(b)| = 0$.

\subsection*{The companion harmonic identity.}

The signal we record alongside the partition statistic is the
squared modulus of the normalized exponential sum at frequency
$1/f$ on the alphabet:
\begin{equation}
  R(k, f) := |S(k, f)|^2,
  \qquad
  S(k, f) := \frac{1}{k} \sum_{j=1}^{k} e^{2\pi i j / f}.
  \label{eq:Rdef}
\end{equation}
For $f \in \Z_{\ge 2}$, $R(k, f)$ admits a closed form via the
geometric sum (Lemma~\ref{lem:countdown} below); the closed form
is the squared Fej\'er kernel at rational frequency, evaluated
at integer abscissa $k$ and normalized by alphabet size $k$. The
closed form holds for every integer $f \ge 2$, not specifically
primes; primality enters only when interpreting $f = p_1 =
\spf(b)$ as the smallest prime factor of an integer $b$.

\subsection*{What this paper records.}

\begin{enumerate}[label={\upshape(\arabic*)}]
\item The Fej\'er-kernel closed form
(Lemma~\ref{lem:countdown}), with attribution to Fej\'er 1900,
Apostol 1976, Iwaniec--Kowalski 2004, and Oppenheim--Schafer
2010 as the standard genealogy.
\item The Eratosthenes synchronization
(Theorem~\ref{thm:zerowidth}): the first-coprime-failure event
at $k = p_1$ and the first integer zero of $R(\cdot, p_1)$ at
$k = p_1$ are two coordinate-translations of the same
elementary fact ``the smallest positive multiple of $p_1$ is
$p_1$.''
\item The $\omega$-blindness corollary
(Corollary~\ref{cor:omegablind}): $R(k, 1/p)$ depends only on
$k$ and $p$, not on the modulus $b$.
\item The continuum identification
(Theorem~\ref{thm:continuum}): $R(k, f) \to \sinc^2(k/f)$ as
$f \to \infty$ along $k/f \to t$ fixed, recovering
$\sinc^2(1/2) = 4/\pi^2$ at the half-period.
\item An exhaustive verification harness:
$712$ distinct algebraic checks across $8$ primes and $187$
semiprimes, max deviation $1.11 \times 10^{-16}$, zero
exceptions, runtime under three minutes. Implementation in
$\sim 200$ lines of Python with \texttt{numpy} +
\texttt{cmath}; runnable from the manuscript folder.
\end{enumerate}

\subsection*{Honest accounting of novelty.}

The closed form (Lemma~\ref{lem:countdown}) is standard. The
synchronization (Theorem~\ref{thm:zerowidth}) is a
coordinate translation of the elementary fact ``$p_1$ is the
smallest positive multiple of $p_1$.'' The $\omega$-blindness
corollary is a one-line consequence of the closed form. The
continuum identity is the standard discrete-to-continuum limit
of the Fej\'er kernel.

The contribution of the present paper is the \emph{packaging}
plus the verification harness: collecting these four facts
together, identifying the synchronization between the arithmetic
sieve and the harmonic kernel, and verifying the synchronization
exhaustively across all primes up to $23$ and all semiprimes
$b = pq$ with $3 \le p < q$, $p \le 59$. The harness is
$712$ algebraic checks at machine precision; we record this as
a clean object of study in the spirit of \emph{Experimental
Mathematics}.

\subsection*{The 712 vs 36{,}662 question.}

The present manuscript reports $712$ distinct algebraic checks,
the harness count for the four specific theorems. A separate
working paper~\cite{Sanders2026WP35} of the broader research
program reports a cumulative $36{,}662$ exact computations
across a wider corpus of related identities and observations
(seeded residue persistence, $D_1/D_2$ kinematics, scale-free
sinc approximations, etc.). The $712$ is what the present
manuscript verifies; the $36{,}662$ is the cumulative corpus
total of which $712$ is the J08-specific subset. We mention the
$36{,}662$ here only to forestall confusion; the manuscript's
verification claims are the $712$ at machine precision, full
stop.

%%% ==================================================================
\section{Setup}
\label{sec:setup}
%%% ==================================================================

Throughout, $b$ denotes an integer with $b > 1$ and smallest prime
factor $p_1 = \spf(b)$. We use the coprimality partition of
\cite{J04Sanders}: for $k \in \{1, 2, \ldots\}$,
\[
  C_k(b) = \{ x \in \{1, \ldots, k\} : \gcd(x, b) = 1 \},
  \qquad
  G_k(b) = \{1, \ldots, k\} \setminus C_k(b).
\]
The First-G Localization Lemma of J04 \cite[Theorem~3.1]{J04Sanders}
asserts $|G_k(b)| = 0$ for every $k < p_1$ and
$G_{p_1}(b) = \{p_1\}$.

\begin{definition}[Harmonic resonance signal]
\label{def:R}
For an integer $f \ge 2$ and $k \ge 1$, define
\begin{equation}
  S(k, f) = \frac{1}{k} \sum_{j=1}^{k} e^{2\pi i j / f},
  \qquad
  R(k, f) = |S(k, f)|^2.
  \label{eq:RSdef}
\end{equation}
\end{definition}

\begin{remark}
$S(k, f)$ is the discrete-time Fourier transform of the indicator
$\mathbf{1}_{\{1, \ldots, k\}}$ evaluated at frequency $1/f$,
normalized by alphabet size. We write $R(k, f)$ throughout, with
$f$ being the prime (or other integer $\ge 2$); some sources use
the equivalent notation $R(k, 1/p)$ for emphasis on the frequency.
\end{remark}

%%% ==================================================================
\section{The discrete Fej\'er kernel at rational frequency}
\label{sec:lemma}
%%% ==================================================================

\begin{lemma}[Discrete Fej\'er kernel; cf.\ Fej\'er 1900,
Apostol 1976 \S11.5, Iwaniec--Kowalski 2004 \S1.7]
\label{lem:countdown}
For every integer $f \ge 2$ and every positive integer $k$,
\begin{equation}
  R(k, f) = \frac{\sin^2(\pi k / f)}{k^2 \sin^2(\pi / f)}.
  \label{eq:countdown}
\end{equation}
Hence
\[
  R(1, f) = 1,
  \quad
  R(f-1, f) = \frac{1}{(f-1)^2},
  \quad
  R(f, f) = 0,
  \quad
  R(f+1, f) = \frac{1}{(f+1)^2},
\]
and $R(\cdot, f)$ is strictly decreasing on $\{1, 2, \ldots, f-1\}$.
The integer zero $R(f, f) = 0$ has zero width: $R(f-1, f) > 0$
and $R(f+1, f) > 0$.
\end{lemma}

\begin{proof}
The geometric-sum formula gives
\[
  \sum_{j=1}^{k} e^{2\pi i j / f}
  = e^{2\pi i / f} \cdot \frac{1 - e^{2\pi i k / f}}{1 - e^{2\pi i / f}}.
\]
Taking the squared modulus of $(1/k)$ times this sum and applying
$|1 - e^{i\theta}|^2 = 4 \sin^2(\theta/2)$ twice gives
\[
  |S(k, f)|^2
  = \frac{1}{k^2} \cdot \frac{|1 - e^{2\pi i k / f}|^2}{|1 - e^{2\pi i / f}|^2}
  = \frac{1}{k^2} \cdot \frac{4 \sin^2(\pi k / f)}{4 \sin^2(\pi / f)}
  = \frac{\sin^2(\pi k / f)}{k^2 \sin^2(\pi / f)},
\]
giving \eqref{eq:countdown}. The boundary values follow by
substitution: at $k = 1$, $R = \sin^2(\pi/f)/\sin^2(\pi/f) = 1$;
at $k = f$, $\sin(\pi) = 0$; at $k = f \pm 1$,
$\sin^2(\pi(f \pm 1)/f) = \sin^2(\pi/f)$ via
$\sin^2(\pi - x) = \sin^2(x)$, so
$R(f \pm 1, f) = 1/(f \pm 1)^2$. Strict decrease on
$\{1, \ldots, f-1\}$ follows because the function
$x \mapsto \sin(\pi x / f) / x$ is strictly decreasing on
$(0, f)$ (its derivative is negative there since
$x \mapsto \sin(x)/x$ is strictly decreasing on $(0, \pi)$); the
strictness propagates to integer abscissae and is preserved by
squaring of non-negative values.
\end{proof}

\begin{remark}[Primality is not used]
The closed form \eqref{eq:countdown} holds for every integer
$f \ge 2$. The discrete Fej\'er kernel's primality-blindness is
what makes the synchronization in
Section~\ref{sec:synchronization} possible: at $f = p_1 =
\spf(b)$, both the arithmetic sieve event ``$|G_k(b)|$ first
nonzero'' and the harmonic event ``$R(k, p_1) = 0$'' occur at
$k = p_1$ for the same elementary reason ($p_1 \mid k$),
expressed in two coordinate systems.
\end{remark}

%%% ==================================================================
\section{Eratosthenes synchronization}
\label{sec:synchronization}
%%% ==================================================================

\begin{theorem}[Coordinate translation between arithmetic and
harmonic gates]
\label{thm:zerowidth}
Let $b > 1$ with smallest prime factor $p_1 = \spf(b)$. The
following two events occur at the same integer $k = p_1$:
\begin{itemize}\setlength\itemsep{2pt}
\item[\rm (i)] (Arithmetic gate.) The smallest $k \ge 1$ for which
$|G_k(b)| > 0$ is $k = p_1$, with $G_{p_1}(b) = \{p_1\}$
\cite[Theorem~3.1]{J04Sanders}.
\item[\rm (ii)] (Harmonic zero.) The smallest $k \ge 1$ for
which $R(k, p_1) = 0$ is $k = p_1$, with $R(p_1 - 1, p_1) =
1/(p_1 - 1)^2 > 0$ (Lemma~\ref{lem:countdown}).
\end{itemize}
Both (i) and (ii) are reformulations of the elementary fact
\emph{the smallest positive multiple of $p_1$ is $p_1$}.
\end{theorem}

\begin{proof}
The smallest positive integer that has $p_1$ as a divisor is
$p_1$ itself. Statement (i) reformulates this as ``$p_1$ is the
smallest $k$ with $\gcd(k, b) > 1$'' (because $p_1 \mid b$ and
no smaller positive integer has any prime factor of $b$).
Statement (ii) reformulates the same fact as ``$p_1$ is the
smallest $k$ with $\sin(\pi k / p_1) = 0$'' (because
$\sin(\pi k / p_1) = 0$ iff $p_1 \mid k$). The two events
co-locate at $k = p_1$ because they encode the same divisibility
condition.
\end{proof}

\begin{remark}[The synchronization is a tautology, recorded for
clarity]\label{rem:tautology}
Theorem~\ref{thm:zerowidth} is not a coincidence between two
independently-defined objects. The arithmetic event ``$|G_k(b)|$
first nonzero'' and the harmonic event
``$R(\cdot, p_1)$ first zero'' both solve ``smallest positive
$k$ with $p_1 \mid k$'' --- the first phrased in terms of
divisibility of $b$ by $p_1$, the second in terms of vanishing
of $\sin(\pi k / p_1)$. We record the synchronization because
it is useful: a reader presented with either signal can compute
$p_1$ from the gate position. But the synchronization itself is
a re-coordinatization, not a non-trivial coincidence requiring
explanation.
\end{remark}

\begin{corollary}[$\omega$-blindness of the harmonic resonance]
\label{cor:omegablind}
For a fixed prime $p$, the function $k \mapsto R(k, 1/p)$ depends
only on $k$ and $p$ and not on the modulus $b$, regardless of the
ring structure $\omega(b)$ (number of distinct prime factors of $b$).
Concretely: if $p \mid b_1$ and $p \mid b_2$ for two integers
$b_1, b_2$, then $R(k, 1/p)$ has the same value at $b = b_1$ and
$b = b_2$ for every $k$.
\end{corollary}

\begin{proof}
By Lemma~\ref{lem:countdown},
$R(k, 1/p) = \sin^2(\pi k / p) / (k^2 \sin^2(\pi/p))$, which is
a function of $k$ and $p$ alone. The modulus $b$ does not appear
in the closed form.
\end{proof}

\begin{remark}[Consequence for ring-structure detection]
\label{rem:omega-detection}
An observer with access only to the harmonic resonance signal
$R(k, 1/p)$ for $k = 1, \ldots, p-1$ cannot distinguish $b = p^2$
from $b = p \cdot q$ from $b = p \cdot q \cdot r$. To detect
$\omega(b)$ from the alphabet, one must observe an additional
signal --- e.g., the closure-defect statistic of
\cite[\S6]{Sanders2026WP35}, which does vary with $\omega(b)$.
The harmonic resonance gives the prime; the closure defect gives
the ring.
\end{remark}

%%% ==================================================================
\section{Continuum identification with $\sinc^2$}
\label{sec:continuum}
%%% ==================================================================

\begin{theorem}[Discrete-to-continuum identity]
\label{thm:continuum}
Fix $t \in (0, 1)$ and let $f \to \infty$ along a sequence of
integers with $k / f \to t$. Then
\begin{equation}
  R(k, f) \longrightarrow \sinc^2(t),
  \qquad \sinc(t) := \frac{\sin(\pi t)}{\pi t}.
  \label{eq:continuum}
\end{equation}
At the half-period $t = 1/2$,
\[
  \lim_{f \to \infty,\, k/f \to 1/2} R(k, f) = \sinc^2(1/2) = \frac{4}{\pi^2}
  = 0.4052847\ldots,
\]
with rate of convergence $\bigO(1/f^2)$.
\end{theorem}

\begin{proof}
By Lemma~\ref{lem:countdown},
$R(k, f) = \sin^2(\pi k / f) / (k^2 \sin^2(\pi / f))$. As
$f \to \infty$ with $k / f \to t$ along a sequence of integers
$k = k(f)$, we have $\sin(\pi k / f) \to \sin(\pi t)$ and
$f \sin(\pi / f) = \pi - \pi^3/(6 f^2) + \bigO(1/f^4) \to \pi$,
so $k \sin(\pi/f) = (k/f) \cdot f \sin(\pi/f) \to t \cdot \pi$.
Hence
\[
  R(k, f)
  = \frac{\sin^2(\pi k / f)}{(k \sin(\pi / f))^2}
  \longrightarrow \frac{\sin^2(\pi t)}{(\pi t)^2}
  = \sinc^2(t).
\]
At $t = 1/2$, $\sinc(1/2) = (\sin(\pi/2)) / (\pi/2) = 2/\pi$, so
$\sinc^2(1/2) = 4/\pi^2$. The error term: at $t = 1/2$ along
$k = \lfloor f / 2 \rfloor$, $|R(k, f) - \sinc^2(1/2)| =
\bigO(1/f^2)$ from the Taylor expansion of $\sin^2(\pi/f)$
around $f = \infty$, since
$f \sin(\pi/f) = \pi(1 - \pi^2/(6 f^2) + \bigO(1/f^4))$ and
the leading correction propagates through the squared
denominator with coefficient $\pi^2/3$.
\end{proof}

\begin{remark}[Alphabet as rectangular spectral window]
\label{rem:rectangular}
The continuum function $\sinc^2(t)$ is the power spectral density
of the indicator $\mathbf{1}_{[0, 1]}$ on $\R$ \cite{Shannon1949,
OppenheimSchafer2010}. The discrete alphabet $\{1, \ldots, k\}$
is a discrete rectangular window of width $k$;
Theorem~\ref{thm:continuum} identifies the discrete spectral
power $R(k, f)$ at frequency $1/f$ with the continuous power
spectral density of the rectangular window in the appropriate
scaling limit. This identification is the standard
discrete-to-continuum limit of the Fej\'er kernel.
\end{remark}

%%% ==================================================================
\section{Verification}
\label{sec:verification}
%%% ==================================================================

We verified Lemma~\ref{lem:countdown},
Theorems \ref{thm:zerowidth} and \ref{thm:continuum}, and
Corollary~\ref{cor:omegablind} exhaustively at machine precision.
The verification harness comprises four specific check classes;
the totals are tabulated in
Section~\ref{sec:verify-totals}. All computations use IEEE
double-precision floating point (errors at machine epsilon
$\sim 10^{-16}$).

\subsection{Closed-form identity at small primes}

For each prime $f \in \{3, 5, 7, 11, 13, 17, 19, 23\}$, we
compute $R(k, f)$ as a literal double-precision sum
$|(1/k) \sum_{j=1}^{k} e^{2\pi i j / f}|^2$ and compare to the
closed form \eqref{eq:countdown} for all $k \in \{1, \ldots, f+1\}$.
Total $(f, k)$ pairs: $106$. Maximum error across the
$106$ comparisons: $4.44 \times 10^{-16}$. Selected pre-collapse
values:
\begin{center}
\begin{tabular}{c|c|c|c|c}
\toprule
$f$ & $k = f - 1$ & $R(f-1, f)$ measured & $1 / (f-1)^2$ predicted & error \\
\midrule
$3$  & $2$  & $0.25000000$  & $0.25000000$  & $5.55 \times 10^{-17}$ \\
$5$  & $4$  & $0.06250000$  & $0.06250000$  & $2.78 \times 10^{-17}$ \\
$7$  & $6$  & $0.02777778$  & $0.02777778$  & $3.47 \times 10^{-18}$ \\
$11$ & $10$ & $0.01000000$  & $0.01000000$  & $1.39 \times 10^{-17}$ \\
$13$ & $12$ & $0.00694444$  & $0.00694444$  & $1.91 \times 10^{-17}$ \\
$17$ & $16$ & $0.00390625$  & $0.00390625$  & $9.54 \times 10^{-18}$ \\
$19$ & $18$ & $0.00308642$  & $0.00308642$  & $1.30 \times 10^{-18}$ \\
$23$ & $22$ & $0.00206612$  & $0.00206612$  & $1.73 \times 10^{-18}$ \\
\bottomrule
\end{tabular}
\end{center}

\subsection{Macro sweep across 187 semiprimes}

For each of $187$ semiprimes $b = p \times q$ with $3 \le p < q$
and $p \le 59$ (and a representative sample of larger $q$), we
verified Lemma~\ref{lem:countdown} at $k = p - 1$ (the
pre-collapse minimum) and Theorem~\ref{thm:zerowidth}
synchronization (the partition gate event at $k = p_1$). The
macro sweep totals $561$ algebraic checks; the maximum error
against the closed form is $1.11 \times 10^{-16}$ (machine
$\epsilon$). Zero exceptions in $561$ cases.

Selected entries:
\begin{center}
\begin{tabular}{c|c|c|c|c|c}
\toprule
$b$ & $p$ & $q$ & $R(p-1, 1/p)$ measured & $1/(p-1)^2$ predicted & error \\
\midrule
$35$   & $5$  & $7$  & $0.062500$ & $0.062500$ & $2.78 \times 10^{-17}$ \\
$77$   & $7$  & $11$ & $0.027778$ & $0.027778$ & $1.39 \times 10^{-17}$ \\
$143$  & $11$ & $13$ & $0.010000$ & $0.010000$ & $1.39 \times 10^{-17}$ \\
$323$  & $17$ & $19$ & $0.003906$ & $0.003906$ & $1.04 \times 10^{-17}$ \\
$667$  & $23$ & $29$ & $0.002066$ & $0.002066$ & $1.73 \times 10^{-18}$ \\
$1073$ & $29$ & $37$ & $0.001276$ & $0.001276$ & $2.39 \times 10^{-18}$ \\
$4187$ & $53$ & $79$ & $0.000370$ & $0.000370$ & $3.96 \times 10^{-18}$ \\
\bottomrule
\end{tabular}
\end{center}

\subsection{$\omega$-blindness across ring structures}

Corollary~\ref{cor:omegablind} predicts $R(k, 1/p)$ is identical
for every $b$ with $p \mid b$, regardless of $\omega(b)$.
Verification holds $p = 7$ fixed and varies $b$ across
$\omega \in \{1, 2, 3\}$:
\begin{center}
\begin{tabular}{c|c|c|c}
\toprule
$b$    & $\omega(b)$ & ring structure         & $R(1, 1/7), R(2, 1/7), \ldots, R(6, 1/7)$ \\
\midrule
$49$   & $1$ & $7^2$        & $1.0000, 0.8117, 0.5610, 0.3156, 0.1299, 0.0278$ \\
$343$  & $1$ & $7^3$        & $1.0000, 0.8117, 0.5610, 0.3156, 0.1299, 0.0278$ \\
$77$   & $2$ & $7 \cdot 11$ & $1.0000, 0.8117, 0.5610, 0.3156, 0.1299, 0.0278$ \\
$91$   & $2$ & $7 \cdot 13$ & $1.0000, 0.8117, 0.5610, 0.3156, 0.1299, 0.0278$ \\
$1001$ & $3$ & $7 \cdot 11 \cdot 13$ & $1.0000, 0.8117, 0.5610, 0.3156, 0.1299, 0.0278$ \\
$1463$ & $3$ & $7 \cdot 11 \cdot 19$ & $1.0000, 0.8117, 0.5610, 0.3156, 0.1299, 0.0278$ \\
\bottomrule
\end{tabular}
\end{center}
The signal is byte-identical across all six rings (up to floating
point, with errors at $10^{-16}$), as required by
Corollary~\ref{cor:omegablind}. With $7$ values of $k$
($k = 1, \ldots, 6$ plus $k = 7$ checking the zero), the harness
contributes $6 \cdot 7 = 42$ check cells.

\subsection{Universal mid-period constant}

Theorem~\ref{thm:continuum} predicts $R(k, p) \to 4/\pi^2$ as
$p \to \infty$ along $k / p \to 1/2$, with rate $\bigO(1/p^2)$.
Numerical verification at proxy primes:
\begin{center}
\begin{tabular}{c|c|c}
\toprule
$p$   & $R(\lfloor p / 2 \rfloor, p)$ & deviation from $4/\pi^2 = 0.4052847\ldots$ \\
\midrule
$1009$    & $0.406090$ & $+8.0 \times 10^{-4}$ \\
$10007$   & $0.405366$ & $+8.1 \times 10^{-5}$ \\
$100003$  & $0.405293$ & $+8.1 \times 10^{-6}$ \\
\bottomrule
\end{tabular}
\end{center}
The successive ratios $8.0/0.81 \approx 100$ and $8.1/0.81
\approx 100$ confirm the predicted $\bigO(1/p^2)$ rate (each
order-of-magnitude increase in $p$ reduces the error by
$\sim 10^2$). Three continuum-limit witnesses; total $3$ check
cells.

\subsection{Verification harness totals}
\label{sec:verify-totals}

The four check classes contribute:
\begin{center}
\begin{tabular}{l|c}
\toprule
Check class & Algebraic checks \\
\midrule
Closed form (Lemma~\ref{lem:countdown}) at 8 primes
& $106$ \\
Synchronization (Theorem~\ref{thm:zerowidth}) at 187 semiprimes
& $561$ \\
$\omega$-blindness (Corollary~\ref{cor:omegablind}) at 6 rings
& $42$ \\
Continuum (Theorem~\ref{thm:continuum}) at 3 large primes
& $3$ \\
\midrule
Total
& $\mathbf{712}$ \\
\bottomrule
\end{tabular}
\end{center}
Maximum deviation across all $712$ checks: $1.11 \times 10^{-16}$
(machine epsilon). Zero exceptions; runtime under three minutes
on a 2024 consumer laptop.

\subsection{Software, runtime, and reproducibility}

The verification script
\texttt{verify\_prime\_phase\_transition.py} runs with
\texttt{numpy}, \texttt{cmath}, and \texttt{math.gcd} on Python
3.10+. Total runtime under three minutes; typically under thirty
seconds for the small-prime closed-form check, with the
$187$-semiprime sweep dominating the runtime. All sums are
computed in double-precision floating point; gcd computations
use exact integer arithmetic. The script exits $0$ if and only
if every comparison falls below $10^{-10}$ (well above machine
epsilon). The bundle DOI is \texttt{10.5281/zenodo.18852047}.

%%% ==================================================================
\section{The constant $4/\pi^2$ in $\sinc^2$ and in
Montgomery's pair correlation: a rectangular-window remark}
\label{sec:montgomery}
%%% ==================================================================

We close with a brief remark on a coincidence of basis functions.
Montgomery~\cite{Montgomery1973} proved (assuming the Riemann
hypothesis, for the relevant range of test parameters) that the
pair correlation of normalized non-trivial zeros of $\zeta(s)$
satisfies
\[
  R_2(u) = 1 - \left(\frac{\sin(\pi u)}{\pi u}\right)^2 = 1 - \sinc^2(u).
\]
Theorem~\ref{thm:continuum} of the present paper, in the
discrete arithmetic side, gives
$\lim R(k, f) = \sinc^2(t)$ as $f \to \infty$ along $k/f \to t$.
At the half-period $t = u = 1/2$, both functions evaluate to
involve the same constant: $\sinc^2(1/2) = 4/\pi^2$ and
$1 - \sinc^2(1/2) = 1 - 4/\pi^2$.

The shared appearance of $\sinc^2$ in the two contexts reflects
the universality of the rectangular spectral window:
$\sinc^2$ is the power spectrum of the indicator of an interval,
which appears in arithmetic (the alphabet $\{1, \ldots, k\}$ in
the present paper) and in spectral statistics (the test function
in Montgomery's pair correlation, where it is the rectangular
window in the explicit formula). This is a common
rectangular-window origin, not a connection between the two
arithmetic phenomena. We do not claim Theorem~\ref{thm:continuum}
bears on the Riemann hypothesis, on the analytic continuation of
$\zeta$, or on any aspect of GUE-type spectral statistics for
zeta zeros.

\begin{remark}[Identity sinc$^2(1/2) = (2/3)/\zeta(2)$]
\label{rem:sinc-zeta}
The half-period value admits the alternative form
\[
  \sinc^2(1/2) = \frac{4}{\pi^2} = \frac{2}{3} \cdot \frac{6}{\pi^2}
  = \frac{2}{3} \cdot \frac{1}{\zeta(2)},
\]
using $\zeta(2) = \pi^2/6$ (Euler 1734). We record this as a
\emph{structural rhyme}, not a theorem of the present paper:
the identity is a one-line algebraic consequence of the
Basel-problem solution and is recognizable to any number theorist
in five seconds. It does not mean that the corridor midpoint of
the rectangular spectral window is forced to $1/2$ by
zeta-function considerations. Why the natural object of study is
the rectangular window $\{1, \ldots, k\}$ rather than some other
weighting is itself an open question. Outside scope.
\end{remark}

\begin{remark}[Open]\label{rem:open-bridge}
The discrete sum $\sum_{k = 1}^{f - 1} R(k, f) \cdot g(k / f)$ for
a smooth test function $g$ converges in the limit $f \to \infty$
to $\int_0^1 \sinc^2(x) g(x)\, dx$ by the Riemann-sum
interpretation of Theorem~\ref{thm:continuum}. Whether the
appropriate dual sum
$\sum_{|\gamma| \le T} R_2(\gamma / 2\pi) \cdot \tilde g(\gamma)$
over Riemann zeros admits an analogous closed-form approximation
in some scaling regime is an open question we leave to future
work; we make no claim here that such an approximation exists.
\end{remark}

%%% ==================================================================
\section{Scope and limitations}
\label{sec:scope}
%%% ==================================================================

We list explicitly what the paper does and does not claim.

\begin{enumerate}[label={\upshape(\arabic*)}]
\item \textbf{What is recorded.} The Fej\'er-kernel closed form
(Lemma~\ref{lem:countdown}, attributed); the
Eratosthenes-vs-harmonic synchronization
(Theorem~\ref{thm:zerowidth}, identified as a coordinate
translation); the $\omega$-blindness corollary
(Corollary~\ref{cor:omegablind}, a one-line consequence of the
closed form); the continuum identity
(Theorem~\ref{thm:continuum}, the standard discrete-to-continuum
limit).

\item \textbf{What is verified.} $712$ distinct algebraic checks
across $8$ small primes, $187$ semiprimes, $6$ ring structures,
and $3$ continuum-limit witnesses. Maximum deviation
$1.11 \times 10^{-16}$ (machine $\epsilon$); zero exceptions;
runtime under three minutes on a 2024 consumer laptop.

\item \textbf{What is \emph{not} claimed.} No statement about
the analytic continuation of $\zeta(s)$. No consequence for the
Riemann hypothesis (the Montgomery-correspondence remark in
Section~\ref{sec:montgomery} is a coincidence of the
rectangular-window basis function, not a mechanism). No
polynomial-time factoring algorithm: Lemma~\ref{lem:countdown}
is pointwise efficient, but extracting $p_1$ from the gate
position $k = p_1$ requires the alphabet to traverse
$\bigO(p_1)$ steps, which is exponential in the bit-length of
an RSA-scale modulus and consistent with sub-exponential
best-known classical complexity (NFS,
\cite{LenstraEtAl1993}). No claim about non-squarefree $b$
beyond what the algebra already covers; the
$\omega$-blindness statement
(Corollary~\ref{cor:omegablind}) is verified for
$\omega(b) \in \{1, 2, 3\}$ and follows formally from the
closed form for any $\omega(b) \ge 1$.

\item \textbf{Restriction to squarefree $b$ in the verification
harness.} The verification script restricts to squarefree $b$
because that is the case shared with the J04 First-G
Localization companion~\cite{J04Sanders}. The closed form
(Lemma~\ref{lem:countdown}) and the $\omega$-blindness
statement (Corollary~\ref{cor:omegablind}) are not restricted
to squarefree $b$; their proofs use only $f \in \Z_{\ge 2}$
and the geometric sum.

\item \textbf{Finite-range verification.} The exhaustive check is
for $p \le 59$ and a representative sample of larger $q$. The
theorems hold in full generality; the finite verification is
reassurance, not a pillar of the result.
\end{enumerate}

%%% ==================================================================
\section*{Acknowledgements}
%%% ==================================================================

This paper is part of a coordinated J-series of submissions by
the authors. The First-G Localization Lemma of the companion
submission~\cite{J04Sanders} is the foundational lemma on which
the synchronization rests. Algorithmic verification was
implemented and checked over the spring of 2026.

%%% ==================================================================
\begin{thebibliography}{99}
%%% ==================================================================

\bibitem{J04Sanders}
B.~R. Sanders, M.~Gish, \emph{The First-G Event in the
Coprimality Partition: Stability Windows, CRT Idempotent Count,
and Prime-Indexed Phase Transitions}, manuscript in preparation,
2026.

\bibitem{FourCore}
B.~R. Sanders, M.~Gish, \emph{Joint Closure, Per-Coordinate Fuse
Data, and a Closed-Form Algebraic Attractor of Two Commutative
Binary Operations on $\Z/10\Z$}, manuscript in preparation, 2026.

\bibitem{Sanders2026WP35}
B.~R. Sanders, M.~Gish, \emph{Prime-Indexed Phase Transitions:
Harmonic Pre-Echo, Closure-Defect, and Scale-Free Sinc
Approximations}, working paper WP35, 7Site LLC research
register, 2026.

\bibitem{Apostol1976}
T.~M. Apostol, \emph{Introduction to Analytic Number Theory},
Undergraduate Texts in Mathematics, Springer-Verlag, New York-Heidelberg,
1976.

\bibitem{Fejer1900}
L.~Fej\'er, \emph{Sur les fonctions born\'ees et int\'egrables},
C.~R.~Acad.~Sci.\ Paris \textbf{131} (1900), 984--987.

\bibitem{IwaniecKowalski2004}
H.~Iwaniec, E.~Kowalski, \emph{Analytic Number Theory}, AMS
Colloquium Publications, Vol.\ 53, American Mathematical Society,
Providence, RI, 2004.

\bibitem{LenstraEtAl1993}
A.~K. Lenstra, H.~W. Lenstra Jr., M.~S. Manasse, J.~M. Pollard,
\emph{The number field sieve}, in \emph{The Development of the
Number Field Sieve}, Lecture Notes in Mathematics 1554,
Springer, Berlin, 1993, pp.~11--42.

\bibitem{Montgomery1973}
H.~L. Montgomery, \emph{The pair correlation of zeros of the zeta
function}, in \emph{Analytic Number Theory} (St.~Louis, MO, 1972),
Proc.~Sympos.~Pure Math., Vol.~XXIV, American Mathematical Society,
Providence, R.I., 1973, pp.~181--193.

\bibitem{Odlyzko1987}
A.~M. Odlyzko, \emph{On the distribution of spacings between zeros
of the zeta function}, Math.\ Comp.\ \textbf{48} (1987), no.~177,
273--308.

\bibitem{OppenheimSchafer2010}
A.~V. Oppenheim, R.~W. Schafer, \emph{Discrete-Time Signal
Processing}, 3rd ed., Prentice Hall, Upper Saddle River, NJ, 2010.

\bibitem{Shannon1949}
C.~E. Shannon, \emph{Communication in the presence of noise},
Proc.~IRE \textbf{37}(1) (1949), 10--21.

\end{thebibliography}

\end{document}
